% ===========================================================================
%  sviluppatore/06-dashboard/05-esportazione.tex
% ===========================================================================
\section{Esportare i risultati}\label{sez:esportazione}

\pulsantesec{Export}, nella barra delle metriche di stream ed ensemble,
scrive in un file \textbf{i dati che la pagina ha disegnato}. Non c'è un
secondo calcolo: sarebbe una seconda versione dei numeri, e prima o poi
diversa.

\subsection{Chi raccoglie e chi scrive}

\begin{center}
\begin{tabular}{@{}L{46mm}L{98mm}@{}}
\toprule
\textbf{Pezzo} & \textbf{Che cosa fa} \\
\midrule
\rigaTenue
\file{static/js/dashboard.js} & Conserva il payload delle API in
  \cod{window.\_\_RISULTATI} ed emette \cod{lyra:risultati}. \\
\file{static/js/metrics\_view.js} & \simbolo{MetricsView.esporta()}: chiavi
  visibili, ordine dei riquadri, metriche personali con i risultati;
  \simbolo{MetricsView.calcola(ids)}: le nascoste, con la stessa richiesta. \\
\rigaTenue
\file{static/js/esporta\_risultati.js} & La finestra delle scelte, le
  immagini (\cod{Plotly.toImage}; le metriche mai disegnate si disegnano fuori
  dallo schermo), la richiesta, lo scaricamento. \\
\file{ui\_ext\_esporta.py} & \rotta{POST}{/api/export/results}: il file, con il
  nome in \cod{Content-Disposition} e in \cod{X-Export-Filename}. \\
\rigaTenue
\file{esporta\_risultati.py} & Dal payload alle \emph{sezioni} (colonne con
  il loro formato, il grafico che le descrive, l'immagine), poi XlsxWriter o
  \pkg{csv}. \\
\bottomrule
\end{tabular}
\end{center}

\subsection{Le regole del file}

\begin{itemize}
  \item Le percentuali delle API arrivano in punti (\num{12.3}); in Excel si
    scrivono come frazioni con il formato percentuale, nei CSV restano in punti
    con \cod{(\%)} nell'intestazione.
  \item Serie con date diverse si allineano per data: un grafico a linee di
    Excel usa le categorie della prima serie per tutte.
  \item Le heatmap diventano tabelle con la scala a tre colori; i pesi, le
    barre nei dati. Le figure con l'asse x numerico, grafici a dispersione.
  \item Nel foglio unico il rapporto è largo quanto la tabella più larga che
    contiene: nessuna sua cella finisce dove stanno le serie lunghe.
\end{itemize}

\begin{pratica}{Aggiungere un riquadro esportabile}
Un riquadro nuovo con \cod{data-sv} si esporta se \simbolo{Costruttore} sa
farne una sezione: un metodo che legge il suo campo del payload e restituisce
una \simbolo{Sezione}, aggiunto all'elenco di \simbolo{Costruttore.sezioni}.
Se ha un grafico Plotly, l'immagine la prende da sé il client.
\end{pratica}

\noindent
La libreria (\pkg{XlsxWriter}) è in \file{requirements.txt} e fra gli
import impliciti di \file{build/lyra-backend.spec}, che ora dichiara anche i
moduli \file{ui\_ext\_*.py}: \file{app.py} li carica per nome e PyInstaller da
solo non li vedrebbe. La prova è \file{tests/prova\_esportazione.py}.
